\hypersetup{pdftitle={\TituloVolumen},pdfauthor={\AutorEdicion},pdfsubject={Análisis. Volumen \NumeroVolumen. Apuntes y ejercicios propuestos},pdfkeywords={\PalabrasClave}}
\makeindex[title={Índice alfabético},intoc,options=-s ../comun/indice.ist]
\setcounter{tocdepth}{1}
\setcounter{secnumdepth}{2}
\setlength{\parindent}{1.2em}
\setlength{\parskip}{3pt plus 1pt minus 1pt}
\emergencystretch=5em
\allowdisplaybreaks[1]
\setlist[enumerate,1]{label=(\roman*),leftmargin=*,itemsep=2pt}
\setlist[itemize]{leftmargin=*,itemsep=2pt}
\pagestyle{fancy}
\fancyhf{}
\fancyhead[L]{\small\nouppercase{\leftmark}}
\fancyfoot[C]{\thepage}
\renewcommand{\headrulewidth}{0.3pt}
\renewcommand{\chaptermark}[1]{\markboth{\thechapter.\ #1}{}}
\fancypagestyle{plain}{\fancyhf{}\fancyfoot[C]{\thepage}\renewcommand{\headrulewidth}{0pt}}
\theoremstyle{plain}
\newtheorem{teorema}{Teorema}[chapter]
\newtheorem{proposicion}[teorema]{Proposición}
\newtheorem{lema}[teorema]{Lema}
\newtheorem{corolario}[teorema]{Corolario}
\theoremstyle{definition}
\newtheorem{definicion}[teorema]{Definición}
\newtheorem{ejemplo}[teorema]{Ejemplo}
\newtheorem{ejercicio}{Ejercicio}[chapter]
\theoremstyle{remark}
\newtheorem{observacion}[teorema]{Observación}
\newenvironment{solucion}{\begin{proof}[Solución]}{\end{proof}}
\newenvironment{estrategia}{\par\noindent\textbf{Estrategia.}\ }{\par}
\newcommand{\R}{\mathbb R}
\newcommand{\C}{\mathbb C}
\newcommand{\N}{\mathbb N}
\newcommand{\Q}{\mathbb Q}
\newcommand{\Z}{\mathbb Z}
\newcommand{\F}{\mathbb F}
\newcommand{\eps}{\varepsilon}
\newcommand{\norm}[1]{\left\lVert#1\right\rVert}
\newcommand{\abs}[1]{\left\lvert#1\right\rvert}
\newcommand{\cl}[1]{\overline{#1}}
\newcommand{\interior}[1]{#1^{\circ}}
\newcommand{\ip}[2]{\left\langle#1,#2\right\rangle}
\newcommand{\Bcal}{\mathcal B}
\newcommand{\Lcal}{\mathcal L}
\newcommand{\Kcal}{\mathcal K}
\DeclareMathOperator{\dist}{dist}
\DeclareMathOperator{\diam}{diam}
\DeclareMathOperator{\spanop}{span}
\DeclareMathOperator{\supp}{supp}
\DeclareMathOperator{\id}{id}
